\section{Introduction}

A missed dose can create a difficult choice: should an older adult manage it independently, or should a family member step in? Family involvement can support medication routines \cite{92}, yet the same checking can shift privacy, responsibility, and control \cite{24,115}. This tension matters especially in many Global South households, where medication work often extends across relatives, shared devices, and intermediated technology use \cite{3,98}. Reminding, purchasing medicines, interpreting prescriptions, and checking doses can therefore be part of everyday family care rather than separate acts of assistance. Medication adherence is more than a problem of remembering. It is also a continuing negotiation over who notices, who acts, and who takes responsibility when a routine breaks down.

Family support can preserve older adults' agency, but it can also redraw its boundaries. Older adults may rely on others for assistance while still wanting control over how support and health information are shared \cite{13,24,30,50}. Interdependence and aging scholarship similarly shows that receiving assistance need not diminish agency because people can rely on others while retaining a meaningful role in how support is arranged \cite{13,30,33}. Yet those boundaries can change across situations: checking may communicate care at one moment and feel like unwanted oversight at another \cite{24,115}. Designing for medication adherence therefore requires support for these shifts without assuming that either independence or family involvement should always take priority.

This negotiation becomes harder when medication technology can act before anyone explicitly asks it to. Mixed-initiative systems have long considered when computational systems should take initiative, while newer proactive agents can intervene in communication or propose actions on a person's behalf \cite{45,52}. In medication care, an AI agent might notice an unconfirmed dose and ask whether a relative should be involved. If the agent keeps the missed dose private, the older adult retains control, but the caregiver may remain unaware. Involving the family can bring support, while also expanding who knows about the dose and who becomes responsible for responding. The AI agent therefore becomes a participant in deciding how responsibility moves between older-adult control and family involvement.

We call this relationship \textbf{allegiance}: whose interests and authority the AI agent gives priority to at a particular moment. Allegiance is not simply a permission setting or a fixed assignment to one user. The same older adult may want private support during an ordinary routine, welcome shared involvement after uncertainty, and resist family intervention in another situation. Because these preferences can shift, allegiance must be visible and revisable rather than hidden in configuration. The design problem is not to align the agent once, but to let people inspect, negotiate, and change whom it serves as circumstances and relationships change.

Care technologies have established that family visibility can support coordination while also redistributing privacy, responsibility, and power. Awareness displays have helped relatives follow older adults' activities and coordinate care \cite{22,82,94}, and collaborative systems have shown how families can participate in person-centred care \cite{81}. Prior studies also show that shared visibility can create new obligations and tensions over who can access information or act on it \cite{24,111,115}. These studies make family participation visible as part of care, not as an external interruption. They leave a further question when the system can itself propose that responsibility move from one person to another.

Research on agent alignment and ways to question automated decisions provides a second foundation for this problem. Alignment research has examined how agents should cooperate with, obey, or represent human objectives \cite{32,38,79}. Work on agents that must account for multiple people's interests further shows how difficult it can be to decide whose interests should take priority \cite{62}. Research on questioning AI decisions, exposing system uncertainty, and revising consent shows how automated actions can remain open to scrutiny and change \cite{6,28,50,71}. Together, these traditions explain why an agent's authority must remain visible and open to revision. What remains unresolved is how that authority is negotiated when an older adult and caregiver both have legitimate, but sometimes different, expectations of what the agent should do.

We study allegiance as an ongoing negotiation rather than a fixed permission or caregiver-access setting. We conducted a two-phase qualitative household study in Bangladesh with 17 older adults and eight family caregivers. Phase 1 examined medication practices and the ways families already shared responsibility. We then developed and deployed a medication-support prototype for two weeks that reminded older adults, recorded medication activity, stated whom it was serving, and requested permission before involving a family member. Post-deployment interviews examined how participants responded to these interactions. Accounts of deployed features were kept separate from judgments about additional mechanisms that were described but not built. Across the study, we analyse when participants treated the agent as a \textbf{tool} for the older adult, a \textbf{coach} shared across the care relationship, or an \textbf{advocate} for greater family involvement. We also examine what made movement among these roles accepted, resisted, or reversed.

Together, these phases trace how responsibility is negotiated before the agent arrives, how allegiance is negotiated after it enters care, and what those negotiations require from design. The study asks:

\textbf{RQ1 (Formative).} How do Bangladeshi older adults and family caregivers distribute and negotiate medication work, and how do they understand responsibility for that work before an agent enters the household?

\textbf{RQ2 (Interaction).} When an agent joins this care network, how do older adults and caregivers negotiate its allegiance, and when are shifts between older-adult control and family involvement accepted, resisted, or reversed?

\textbf{RQ3 (Design).} What makes the agent's roles as tool, coach, or advocate legitimate to older adults and caregivers, and what design implications follow for making shifts among these roles visible, revocable, negotiable, and dignity-preserving?

This work makes three contributions:

\begin{itemize}
  \item \textbf{Empirical.} This study shows how Bangladeshi older adults and family caregivers negotiate medication responsibility and family involvement across formative and post-deployment accounts.
  \item \textbf{Methodological.} We use a caregiver-inclusive household approach that treats caregivers as participants in their own right and preserves divergent accounts of shared medication work.
  \item \textbf{Design.} The study identifies directions for shared-care agents that keep shifts in allegiance visible, negotiable, and reversible.
\end{itemize}

For shared-care agents, whom the system serves cannot remain a hidden configuration choice.
